\subsection{Diagonalizability and Minimal Polynomial Tests}
\begin{theorem} \label{thm:diagonalizable-iff-minimal-polynomial-squarefree} \hlt{Diagonalizability and the Minimal Polynomial}\\
Let $A\in\operatorname{Mat}(n,n,\mathbb{C})$ have distinct eigenvalues $\lambda_1,\ldots,\lambda_d$.\\
Then $A$ is diagonalizable if and only if $\mu_A(z)=(z-\lambda_1)\cdots(z-\lambda_d)$.\\
Equivalently, $A$ is diagonalizable if and only if $\mu_A$ is square-free.
\begin{proof}
Set $p(z)\defeq\prod_{i=1}^{d}(z-\lambda_i)$.
\begin{enumerate}[label=\roman*.]
\setlength{\itemsep}{0pt}
\item[$(\Rightarrow)$] If $A=S\Lambda S^{-1}$ is diagonalizable, then $p(\Lambda)=0$, so $p(A)=0$.\\
Since $\mu_A$ contains every factor $z-\lambda_i$ and divides $p$, we have $\mu_A=p$.
\item[$(\Leftarrow)$] Suppose $\mu_A=p$. On the primary summand $\mathcal{G}_{\lambda_i}(A)$, the restriction of $A-\lambda_jI_n$ for $j\neq i$ is
$(A-\lambda_iI_n)+(\lambda_i-\lambda_j)I$, a nonzero scalar plus a nilpotent operator. By Lemma \ref{lem:scalar-plus-nilpotent-invertible}, these restricted factors are invertible. Since $0=p(A)|_{\mathcal{G}_{\lambda_i}(A)} =\left(\prod\limits_{j\neq i}(A-\lambda_jI_n)\right)|_{\mathcal{G}_{\lambda_i}(A)} (A-\lambda_iI_n)|_{\mathcal{G}_{\lambda_i}(A)}$, the invertible product of the $j\neq i$ factors can be cancelled, giving $(A-\lambda_iI_n)=0$ on $\mathcal{G}_{\lambda_i}(A)$.\\
Thus every primary summand is an eigenspace, so $A$ is diagonalizable.
\end{enumerate}
\hlt{Structural Commentary}: Repeated factors of $\mu_A$ record nilpotent parts in primary blocks.\\
\hlt{Key Insight}: Square-free $\mu_A$ kills each primary block in one step.
\end{proof}
\end{theorem}

\begin{corollary} \label{cor:squarefree-annihilator-implies-diagonalizable} \hlt{Square-Free Annihilator Criterion}\\
Let $A\in\operatorname{Mat}(n,n,\mathbb{C})$. If a polynomial $f$ with no repeated roots satisfies $f(A)=0$, then $A$ is diagonalizable.
\begin{proof}
Since $\mu_A\mid f$, the minimal polynomial is square-free; Theorem \ref{thm:diagonalizable-iff-minimal-polynomial-squarefree} applies.\\
\hlt{Structural Commentary}: A square-free annihilator forces square-free $\mu_A$.\\
\hlt{Key Insight}: Square-free $\mu_A$ means diagonalizable.
\end{proof}
\end{corollary}

\begin{method} \label{meth:operator-minimal-polynomial-from-matrix} \hlt{Computing $\mu_T$ from a Matrix}\\
Let $A\in\operatorname{Mat}(n,n,\mathbb{F})$ be the matrix of $T$ in some basis.
\begin{enumerate}[label=(\roman*)]
\setlength{\itemsep}{0pt}
\item For $m=1,2,\ldots$, test the monic candidate $p_m(z)=a_0+a_1z+\cdots+a_{m-1}z^{m-1}+z^m$.
\item Solve linear system by comparing matrix entries in $a_0I+a_1A+\cdots+a_{m-1}A^{m-1}=-A^m$ for $a_0,\ldots,a_{m-1}$.
\item If no solution exists, increase $m$. The first solvable candidate is $\mu_T$.
\end{enumerate}
\end{method}

\begin{figure}[H]
  \centering
  \begin{tikzpicture}[
  >=Stealth,
  every node/.style={font=\small},
  box/.style={draw=gray!60, fill=gray!6, rounded corners=2pt, align=center, inner xsep=5pt, inner ysep=4.5pt, minimum height=0.74cm, text width=1.9cm},
  main/.style={draw=blue!60!black, fill=blue!6, rounded corners=2pt, align=center, inner xsep=5pt, inner ysep=4.5pt, minimum height=0.74cm, text width=1.9cm},
  good/.style={draw=dartmouthgreen!70!black, fill=dartmouthgreen!8, rounded corners=2pt, align=center, inner xsep=5pt, inner ysep=4.5pt, minimum height=0.74cm, text width=1.5cm},
  arr/.style={->, line width=0.7pt, draw=black!70}
]
  \node[box, text width=2.05cm] (candidate) at (0,0)
    {monic candidate\\$p_m(z)$};
  \node[main] (system) at (2.82,0)
    {solve for\\$a_0,\ldots,a_{m-1}$};
  \node[box] (first) at (5.63,0)
    {first solvable\\degree $m$};
  \node[good] (mu) at (8.00,0)
    {$\mu_T=p_m$};

  \draw[arr] (candidate) -- (system);
  \draw[arr] (system) -- (first);
  \draw[arr] (first) -- (mu);
\end{tikzpicture}

  \caption{Solve $a_0I+\cdots+a_{m-1}A^{m-1}=-A^m$; the first solvable monic candidate is $\mu_T$.}
  \label{fig:minimal-polynomial-linear-system}
\end{figure}

\begin{method} \label{meth:minimal-polynomial-testing-from-characteristic-polynomial} \hlt{Testing a Minimal Polynomial from $p_A$}\\
Suppose $p_A(z)=\prod_i(z-\lambda_i)^{a_i}$ is known.
\begin{enumerate}[label=\roman*.]
\setlength{\itemsep}{0pt}
\item Write each candidate as $m_{\mathbf{q}}(z)=\prod_i(z-\lambda_i)^{q_i}$, where $1\leq q_i\leq a_i$.
\item Test exponent vectors $\mathbf{q}=(q_i)$ in increasing total degree $\sum_iq_i$.
\item For each $\mathbf{q}$, compute $m_{\mathbf{q}}(A)=\prod_i(A-\lambda_iI)^{q_i}$.
\item The first candidate with $m_{\mathbf{q}}(A)=0$ is $\mu_A$, by Theorem \ref{thm:minimal-polynomial-basic-properties}.
\end{enumerate}
\end{method}
